\newcommand{\nomedoc}{Norme di progetto}
\newcommand{\versione}{0.9}
\newcommand{\versioneglossario}{0.1}
\newcommand{\versionenormeprogetto}{0.9}
\newcommand{\nomefile}{NormeDiProgetto-\versione.pdf}
\newcommand{\datacreazione}{7 Dicembre 2010}
\newcommand{\datamodifica}{11 Dicembre 2010}
\newcommand{\stato}{formale}
\newcommand{\uso}{interno}
\newcommand{\redazione}{Baron Federico, Palazzin Alberto}
\newcommand{\verifica}{}
\newcommand{\approvazione}{}
\newcommand{\distribuzione}{
VT.G \\
& Prof. Vardanega Tullio\\
& Prof. Cardin Riccardo }

% FUNZIONI TIPOGRAFICHE
\newcommand{\co}{\texttt} % courier
\newcommand{\bo}{\textbf} % bold
\newcommand{\pr}{\par\medskip} % paragrafo spaziato
\newcommand{\sca}{\textsc} % small caps

\documentclass[a4paper,12pt]{report}
% 10pt,11pt,12pt
% titlepage, notitlepage -> per dare inizio o no ad una nuova pagina dopo titolo
% twoside -> per dire se fronte-retro
\usepackage[latin1]{inputenc}
% per caratteri accentati
\usepackage[italian]{babel}
% per regole sintattiche italiane
\usepackage[bookmarks=true, pdfborder={0 0 0 0}]{hyperref}
% per collegamenti ipertestuali
\usepackage{graphicx}
% per inserimento immagini

% \usepackage{enumerate}
% per personalizzare elenchi puntati

\usepackage[hmargin=2cm]{geometry} %margine 2 cm
%\geometry{options varie}

% comandi per gestire meglio header e footer
\usepackage{fancyhdr}  % header e footer
\usepackage{totpages}
\pagestyle{fancy}
\renewcommand{\headrulewidth}{0.4pt}
\renewcommand{\footrulewidth}{0.4pt}

\setlength{\headheight}{1.2cm} % NON TOCCARE
\setlength{\voffset}{-1.5cm} % NON TOCCARE
\setlength{\textheight}{666pt} % NON TOCCARE
\setlength{\footskip}{60pt}
\setlength{\parindent}{0pt} % INDENTAZIONE

\lhead{\nomedoc\  (ver. \versione)}
\chead{}
\rhead{\includegraphics[height=1cm]{img/netmus.png}}
\lfoot{\includegraphics[height=0.8cm]{img/logo.png}}
\cfoot{}
\rfoot{\thepage}

\usepackage{titlesec}
\titleformat{\chapter}{\normalfont\huge\bfseries}
{\thechapter}{20pt}{\Huge}

\usepackage{rotating}   % PER TABELLE E AMBIENTI RUOTATI
\usepackage{array}
\usepackage{color}
\usepackage{colortbl}  % VARIE PER GESTIRE I COLORI
\definecolor{Orange}{RGB}{255,127,0}   % ARANCIO ACCES0
\definecolor{orange}{RGB}{255,207,80}  % ARANCIO TENUE

\addtocontents{toc}{\protect\thispagestyle{fancy}}  % PER INDICI CON + PAGINE
\usepackage[font=it]{caption}    % PER RENDERE CORSIVE LE DIDASCALIE
\usepackage{eurosym}  % PER SIMBOLO EURO

% \usepackage{listings}   per codice sorgente

\author{VT.G - Valter Texas Group}

\begin{document}

\pagenumbering{Roman} % INIZIO NUMERAZIONE ARABA

\vspace*{1cm}
\begin{center}

\begin{LARGE} \sca{Federico Baron} \end{LARGE}\\
\vspace{0.5cm}
\begin{Large}
\emph{fede.baron.89@gmail.com} \end{Large}\\
\vspace*{1cm} \includegraphics[width=5cm]{img/logo.png}\\
\vspace{0.5cm}
\begin{Large} \emph{``Comunicazione Aumentata/Alternativa per Giovani Ospiti
della Terapia Intensiva Pediatrica''} \end{Large}\\
\vspace{3cm}
\begin{Large} \sca{\nomedoc} \end{Large}\\
\end{center}
\vspace{1cm}

% INFORMAZIONI DOCUMENTO
\begin{center}
\begin{tabular}{r|l}
\hline & \\
\bo{Nome} & \nomefile \\
\bo{Versione attuale} & \versione \\
\bo{Data creazione} & \datacreazione \\
\bo{Data ultima modifica} & \datamodifica \\
\bo{Redazione} & \redazione \\
& \\\hline
\end{tabular}
\end{center}
\newpage

% REGISTRO MODIFICHE
\section*{Registro delle modifiche}
\begin{tabular}{lll}
% REVISIONI DEL DOCUMENTO
% DALLA PIU' NUOVA ALLA PIU' VECCHIA

\bo{Data:} 11/12/2010 &
\bo{Versione:} 0.9 &
\bo{Autore:} Palazzin Alberto\\
\hline\\
\multicolumn{3}{p{470px}}{ Controllo e correzione degli errori ortografici.}\\
\\

\bo{Data:} 10/12/2010 &
\bo{Versione:} 0.8 &
\bo{Autore:} Palazzin Alberto\\
\hline\\
\multicolumn{3}{p{470px}}{ Sistemazione dei Capitoli 5, 6, 7 e 8.}\\
\\

\bo{Data:} 10/12/2010 &
\bo{Versione:} 0.7 &
\bo{Autore:} Palazzin Alberto\\
\hline\\
\multicolumn{3}{p{470px}}{ Sistemazione dei Capitoli 4, 9 e 10.}\\
\\

\bo{Data:} 10/12/2010 &
\bo{Versione:} 0.6 &
\bo{Autore:} Baron Federico\\
\hline\\
\multicolumn{3}{p{470px}}{ Aggiunto il capitolo ``Tracciamento dei
requisiti''.}\\
\\

\bo{Data:} 10/12/2010 &
\bo{Versione:} 0.5 &
\bo{Autore:} Palazzin Alberto\\
\hline\\
\multicolumn{3}{p{470px}}{ Sistemazione del Capitolo 3.}\\
\\

\bo{Data:} 10/12/2010 &
\bo{Versione:} 0.4 &
\bo{Autore:} Palazzin Alberto\\
\hline\\
\multicolumn{3}{p{470px}}{ Aggiunta Sommario e correzione del Capitolo 2.}\\
\\

\bo{Data:} 09/12/2010 &
\bo{Versione:} 0.3 &
\bo{Autore:} Baron Federico\\
\hline\\
\multicolumn{3}{p{470px}}{ Aggiunta nella sezione 2.3 la citazione di nomi di file relativi ad altri documenti.}\\
\\

\bo{Data:} 09/12/2010 &
\bo{Versione:} 0.2 &
\bo{Autore:} Baron Federico\\
\hline\\
\multicolumn{3}{p{470px}}{ Modifica alla sezione 4.3 e correzione accenti.}\\
\\

\bo{Data:} 07/12/2010 &
\bo{Versione:} 0.1 &
\bo{Autore:} Baron Federico\\
\hline\\
\multicolumn{3}{p{470px}}{ Prima stesura del documento Norme Di Progetto.}\\
\\

% TEMPLATE PER REVISIONE

\end{tabular}

% INDICE
\tableofcontents
\thispagestyle{fancy} % per lo stile di header e footer


\chapter*{Sommario}
\thispagestyle{fancy}
Questo documento contiene un elenco di norme, approvate da
tutti i membri del gruppo VT.G e applicate su l'intero progetto; riguardano in
particolare le convenzioni di comunicazione interne ed esterne, riunioni tra i componenti del team e
documentazioni del progetto. 
Inoltre sono definiti alcuni aspetti dell'amministrazione di progetto
particolarmente importanti per poter lavorare con strumenti e tecnologie
comuni.

\thispagestyle{fancy} % serve perche' nelle pagine di inizio Chapter esca header e footer
\pagenumbering{arabic} % INIZIO NUMERAZIONE NORMALE
\rfoot{\thepage\ di \pageref{TotPages}}
\addcontentsline{toc}{chapter}{Sommario}


\chapter{Introduzione}
\thispagestyle{fancy} % serve perche' nelle pagine di inizio Chapter esca header e footer

\section{Scopo del documento}
Il presente documento viene stilato al fine di descrivere le norme che il gruppo
deve seguire per regolare le proprie attivit\`a di lavoro.


\section{Scopo del prodotto}
Il progetto \underline{NetMus} nasce con lo scopo di realizzare un sistema
software basato su \underline{cloud} \underline{computing}, per memorizzare
informazioni di brani musicali in profili utente online.\\ Tali informazioni vengono estratte da
dispositivi musicali o di archiviazione \underline{USB} al momento della loro connessione.

\section{Glossario}
Il Glossario \`e definito con un documento a parte
(\emph{Glossario-\versioneglossario.pdf}). Tutti i termini caratterizzati da
\underline{questa sottolineatura} sono ivi definiti.\\
Verr\`a sottolineata solamente la prima occorrenza di ciascun
termine presente nel Glossario, per non compromettere la leggibilit\`a del documento.

\section{Riferimenti}

\subsection{Normativi} % oppure rif. a Norme di progetto con leggi e tutto
\begin{itemize}
  \item ISO/IEC 12207:1995 - Cicli di vita software
  \item ISO/IEC 9126:2001 - Quality Model
  \item \emph{NormeDiProgetto-\versionenormeprogetto.pdf} che regola e
  accompagna tutti i documenti ufficiali.
\end{itemize}
\newpage
\subsection{Informativi}
\begin{itemize}
  \item Capitolato d'appalto CO2-NETMUS del corso di Ingegneria del Software
  A.A. 2010/11 :\\
  \url{http://www.math.unipd.it/~tullio/IS-1/2010/Progetto/NetMus.pdf}
  \item Slide delle lezioni del corso:\\
  \url{http://www.math.unipd.it/~tullio/IS-1/2010/}
  \item Verbale intervista proponente:\\
  \co{allegato Verbale-1.0.pdf}
  \item Sistema di cloud Google App Engine:\\
  \url{http://code.google.com/intl/it/appengine/}
\end{itemize}


\chapter{Comunicazioni}
\thispagestyle{fancy}
Le comunicazioni si dividono tra interne ed esterne al gruppo, e vengono gestite
nei modi di seguito riportati.

\section{Interne}
Le comunicazioni interne tra i membri del team si svolgono attraverso gli
strumenti offerti da \underline{Google Groups} nel profilo Valter Texas Group e
le e-mail personali. Nello specifico si utilizzano un forum per le discussioni
(sezione ``Discussioni''), dei documenti condivisi \underline{Google Docs} per
alcune note (sezione ``Documenti'') e degli avvisi tramite e-mail. 
Queste comunicazioni sono protette e visibili soltanto ai membri del gruppo.


\section{Esterne}
Le comunicazioni esterne sono gestite con una e-mail di gruppo
valtertexasgroup@googlegroups.com, fornita da Google Groups. I messaggi ricevuti
a questo indirizzo di posta elettronica sono inoltrati automaticamente alle e-mail personali dei componenti del team e pubblicati nella
sezione ``Discussioni'' del forum interno.


\chapter{Norme per la stesura dei documenti}
\thispagestyle{fancy}

\section{Nomenclatura}
Tutti i nomi dei file devono cominciare con la lettera maiuscola ed
indicare in maniera univoca il loro contenuto. Il nome deve essere seguito da un
trattino (-) e due cifre separate da un punto che indicano la versione del
documento (esempio: NormeDiProgetto-1.0.pdf).

\section{Contenuto}
Tutti i documenti devono contenere: 
\begin{itemize}
\item {una pagina iniziale con il logo e il nome del gruppo;}
\item {la sede principale del gruppo;}
\item {il contatto del gruppo;}
\item {il nome del documento;}
\item {la versione del documento;}
\item {la data di creazione;}
\item {la data di ultima modifica;}
\item {lo stato del documento (formale o informale);}
\item {l'uso del documento (esterno o interno);}
\item {la redazione;}
\item {la verifica;}
\item {l'approvazione;}
\item {la distribuzione;}
\item {un registro delle modifiche;}
\item {un indice del documento;}
\item {un sommario contenente un breve riassunto del documento;}
\item {l'indice delle tabelle (se presenta almeno una tabella nel documento);}
\item {l'indice delle figure (se presente almeno una figura nel documento).}
\end{itemize} 

\section{Norme tipografiche}
Il font da utilizzare \`e Knuth's Computer Modern (di default in
\underline{Latex}), la dimensione del capitolo e del sottocapitolo \`e quella di
default in Latex, la dimensione del testo \`e 12pt e la codifica per i caratteri
\`e ISO-8859-1. In ogni documento si deve utilizzare la \underline{sottolineatura} per marcare i termini presenti nel glossario (solo
alla prima ricorrenza) e l'\emph{italic} per evidenziare i nomi di file che fanno riferimento ad altri documenti.

\section{Norme di stile}
Ogni capitolo deve essere identificato da un numero mentre ogni sottocapitolo
deve essere identificato da due numeri separati da un punto (il primo indica il
capitolo e il secondo il sottocapitolo). Tra un capitolo/sottocapitolo e
l'altro, si deve utilizzare un'interlinea come separatore.

\section{Grafici ed illustrazioni}
Qualsiasi immagine e grafico devono essere accompagnati da un numero progressivo
all'interno del documento, preceduto da numeri che ne identificano il
capitolo e il sottocapitolo in cui si trova (esempio: immagine 1.3.1 \`e
l'immagine 1 all'interno del sottocapitolo 3 del primo capitolo) e da una
didascalia che ne dia una breve spiegazione.

\section{Acronimi ed abbreviazioni}
Le diverse revisioni sono indicate con:
\begin{itemize}
\item {RR: Revisione dei Requisiti;}
\item {RPS: Revisione di Progettazione con Specifica Tecnica; }
\item {RPD: Revisione di Progettazione con Defizione del Prodotto;}
\item {RQ: Revisione di Qualifica;}
\item {RA: Revisione di Accettazione.}
\end{itemize}
Si utilizza l'acronimo VT.G (Valter Texas Group) per qualsiasi riferimento al
nostro gruppo di progetto.


\chapter{Strumenti e tecnologie da utilizzare}
\thispagestyle{fancy}
Si consiglia a tutti i membri di VT.G di utilizzare gli stessi strumenti e
configurazioni per qualsiasi attivit\`a di lavoro. Qui sotto ne viene
riportata una lista.

\section{IDE}
Eclipse for Java Developers (versione $\geq$ Helios 3.6).
\\
Scarica e installa dal link
\url{http://www.eclipse.org/downloads/packages/eclipse-ide-java-developers/heliossr1}.
 
\section{Repository}
\underline{Google Code} Project Hosting.
\\
All'interno del progetto (\url{http://code.google.com/p/netmus/}) vengono
utilizzati due \underline{repository}: ``default'' per il progetto
(\url{https://netmus.googlecode.com/hg/}) e ``documents'' per i documenti
(\url{https://documents.netmus.googlecode.com/hg}). Per facilitare il lavoro
simultaneo tra i membri, nel repository documents  sono stati aperti tanti
\underline{branch} quanti sono i documenti da redarre che in prossimit\`a delle
revisioni vengono fatti convergerge.

\section{Sistema di versionamento}
\underline{Mercurial}  (versione $\geq$ 1.7.2). 
\\
Scarica e installa dal link
\url{http://mercurial.selenic.com/}.

\section{Plugin Mercurial per Eclipse}
MercurialEclipse (versione $\geq$ 1.7). 
\\
Per installare seguire queste
operazioni: 
\begin{itemize}
\item {aprire il programma Eclipse;} 
\item {nel men\`u in alto selezionare ``Help'';}
\item {selezionare ``Install new Software'';}
\item {in ``Work with'' inserire \url{http://cbes.javaforge.com/update};}
\item {selezionare ``Mercurial Eclipse'' dalla lista comparsa e deselezionare
``Windows Binaries for Mercurial'' se non si usa un ambiente Windows (se
richiede una login e una password, andare al link
\url{http://www.javaforge.com/createUser.spr}, registrarsi al sito e poi
utilizzare queste credenziali per continuare l'istallazione del plugin);}
\item {installare.}
\end{itemize}

% \section{SDK per AppEngine}
% ???
% 
% \section{Plugin per AppEngine}
% ???

\section{Latex}
\begin{itemize}
\item {Per Linux: TeX Live (versione 2010). 
\\
Scarica e installa dal link
\url{http://www.tug.org/texlive/}.}
\item {Per Mac OS X: MacTex (versione 2010). 
\\
Scarica e installa dal link
\url{http://www.tug.org/mactex/}.}
\item {Per Windows: proTeXt (versione 2010). 
\\
Scarica e installa dal link
\url{http://www.tug.org/protext/}.}
\end{itemize}

\section{Plugin Latex per Ecplise}
TeXlipse (versione $\geq$ 1.4.0). 
\\
Per installare seguire queste
operazioni: 
\begin{itemize}
\item {aprire il programma Eclipse;} 
\item {nel men\`u in alto selezionare ``Help'';}
\item {selezionare ``Install new Software'';}
\item {in ``Work with'' inserire \url{http://texlipse.sourceforge.net};}
\item {selezionare l'unico file dalla lista comparsa;}
\item {installare.}
\end{itemize} 

\section{UML}
ArgoUML (versione $\geq$ 0.30.2) configurato per \underline{UML} 2.x. 
\\
Scarica e installa dal link
\url{http://argouml-downloads.tigris.org/argouml-0.30.2/}. Per eseguirlo utilizzare ``argouml2.sh'' (o ``.exe''), file presente dentro la
directory del programma.

\section{Gantt}
GanttProject (versione $\geq$ 2.0.10). 
\\
Scarica e installa dal link
\url{http://www.ganttproject.biz/download}.


\chapter{Repository e Versionamento}
\thispagestyle{fancy}
Come servizio di hosting per mantenere online i repository
viene utilizzato Google Code Project Hosting
(\url{http://code.google.com/hosting/}). Questo offre un buon
servizio di \underline{ticketing} e una comoda e completa interfaccia web per
la visualizzazione e gestione dei dati caricati. In tale spazio sono
stati creati due repository: uno per i documenti e uno per il codice sorgente
(per maggiori dettagli vedi sezione 4.2). 
\\
Come sistema di
\underline{versionamento} per tali repository viene utilizzato Mercurial
(versione $\geq$ 1.7.2), strumento che deve essere installato nel sistema di
ogni componente del gruppo (per maggiori dettagli vedi sezione 4.3). 
\\
Poich\'e durante l'intero sviluppo verr\`a usato principalmente l'ambiente
Eclipse, si utilizzer\`a il comodo plugin Mercurial Eclipse che consente di accedere
graficamente a tutte le funzionalit\`a di Mercurial (per maggiori dettagli vedi
sezione 4.4). 
\\
Per scaricare localmente un repository con le ultime versioni dei
file si deve eseguire un ``\underline{Pull}'' all'interno del progetto
selezionato. Una volta effettuate delle modifiche, aggiunte o rimozioni viene
utilizzato il relativo ``\underline{Commit}'' per descrivere l'avanzamento di
versione e poi viene ultimato il ``\underline{Push}'' per caricare tale versione
nel repository. Mercurial con i ``\underline{Merge}'' riesce a gestire in maniera efficiente i vari conflitti
che potrebbero esserci in ``Push'' simultanei nello stesso branch. 
\\
In ogni ``Commit'' viene adottata tale convenzione: 
``Operazione NomeDelFile Descrizione''
\begin{itemize}
\item {Operazione: \`e l'operazione fatta su un determinato file (ADD, NEW, FIX,
DEL) e nello specifico}
\begin {itemize}
\item {ADD: \`e stato aggiunto qualcosa di nuovo ad un file gi\`a esistente;}
\item {NEW: \`e stato creato qualcosa di nuovo, quindi un nuovo file;} 
\item {FIX: \`e stato modificato un contenuto presente in un file gi\`a
esistente;} 
\item {DEL: \`e stato totalmente eliminato qualcosa;}
\end {itemize}
\item {NomeDelFile: \`e il nome del file in questione;}
\item {Descrione: \`e una breve descrizione dell'operazione eseguita.}
\end{itemize}
Riguardo al versionamento di un documento, ogni modifica apportata a questo
determina un incremento dell'indice minore (esempio: da 0.1 a 0.2) mentre ogni
rilascio pubblico determina un incremento dell'indice maggiore (esempio in
previsione della RR: da 0.19 a 1.0).


\chapter{Registro delle modifiche}
\thispagestyle{fancy}
Per agevolare la comprensione e tracciare uno storico delle modifiche
effettuate, all'interno di ciascun documento deve essere sempre
presente un registro dettagliato delle modifiche (a partire da
pagina II prima dell'indice). Questo permette di evidenziare la data
d'introduzione della modifica, il numero di versione raggiunto e l'autore di tale modifica.

\chapter{Tracciamento dei requisiti}
\thispagestyle{fancy}
Per identificare in modo univoco e significativo ognuno dei requisiti software
estratti dall'analisi fatta dai VT.G, gli viene assegnato un codice. La lista
dei requisiti sar\`a presente solamente nel documento
\emph{AnalisiDeiRequisiti-1.0.pdf} mentre in tutti gli altri verranno riferiti tramite il codice suddetto.
\\
\\
(C1/C2)(F/Q/I/V)(N/D/O)-n1.n2.n3...nx\\
C1 = componente 1 (componente di persistenza e visualizzazione)\\
C2 = componente 2 (componente di recupero delle informazioni dei brani
musicali contenuti nei sistemi di riproduzione personale dell'utente)\\
F = funzionale\\
Q = qualit�\\
I = interfacciamento\\
V = vincolo\\
N = obbligatorio (necessario)\\
D = desiderabile\\
O = opzionale\\
nx = livello di innestamento nella gerarchia dei requisiti. I ``figli'' di un
requisito avranno una cifra in pi\`u del genitore.
\\

\chapter{Ticketing e time-tracking}
\thispagestyle{fancy}
Il tracciamento delle ore di lavoro per ogni componente verr\`a eseguito tramite
un progetto di GanttProject appositamente condiviso ed accessibile da ognuno
(non concorrentemente). Questo permetter\`a anche un comodo confronto con la
gestione delle risorse pianificata nel PQ.
Il tiketing sar\`a utilizzato tramite gli appositi strumenti forniti da Google Code.



\chapter{Riunioni ed incontri}
\thispagestyle{fancy}
Tutte le riunioni e gli incontri ufficiali vengono svolti nelle seguenti sedi:

\begin{itemize}
\item Residenza Messori, in via Manin 31 a Padova (sede principale);
\item Laboratorio di Informatica, in via Paolotti a Padova.
\end{itemize}


\chapter{Consegna}
\thispagestyle{fancy}
Per ciascuna revisione, il gruppo VT.G si impegna a consegnare i documenti
tramite l'e-mail del gruppo in un unico file compresso ``.zip''. Al suo interno i
file saranno suddivisi in base all'uso, interno o esterno.


\end{document}